\originalchapter{8.701 作业四:弱衰变与 DIS}
题源: Fall 2020 Problem Set 4, 2020年10月21日. 官方答案的教材页码仅作来源索引, 以下推导自足, 不收录第三方教材原文.
\begin{exercise}\label{ex:8701-4-1}
原题 1. 对 $\mu\to e\nu_\mu\bar\nu_e$ 描述计算 $\tau_\mu=192\pi^3/(G_F^2m_\mu^5)$ 所需步骤及假设; 加分要求完整计算.
\end{exercise}
\begin{solution}
官方框架重编, 积分步骤为编者补写. 取 $\mu^-$ 静止、初态不偏振, 忽略 $m_e$ 与中微子质量, $q^2\ll M_W^2$, 不含辐射修正. Fermi 振幅自旋平均后为
$\overline{|M|^2}=64G_F^2(P\cdot p_{\bar\nu_e})(p_e\cdot p_{\nu_\mu})$.
先固定电子动量, 剩两中微子总动量 $Q=P-p_e$. 在 Q 静止系积分角度, 利用 $\int\dd\Phi_2=1/(8\pi)$ 和 $\langle n_in_j\rangle=\delta_{ij}/3$, 可得
\begin{equation}
 \int\dd\Phi_2\,p_{\bar\nu_e}^{\alpha}p_{\nu_\mu}^{\beta}
 =\frac1{96\pi}(Q^2g^{\alpha\beta}+2Q^\alpha Q^\beta).
\end{equation}
同 $P_\alpha p_{e\beta}$ 收缩, 得 $m_\mu^2E_e(3m_\mu-4E_e)/(96\pi)$. 外层电子测度积分角度后为 $E_e\dd E_e/(4\pi^2)$, 再乘 $1/(2m_\mu)$ 得
\begin{equation}
 \frac{\dd\Gamma}{\dd E_e}=\frac{G_F^2m_\mu}{12\pi^3}E_e^2(3m_\mu-4E_e),\quad0<E_e<m_\mu/2.
\end{equation}
$\int_0^{m/2}E^2(3m-4E)\dd E=m^4/16$, 因此 $\Gamma=G_F^2m^5/(192\pi^3)$. 自然单位中 $\tau=1/\Gamma$, 秒单位中 $\tau=\hbar/\Gamma$. 三体相空间、初态平均和 V--A 矩阵元都不可遗漏.
\end{solution}
\begin{exercise}\label{ex:8701-4-2}
原题 2. 比较中子 $n\to p e\nu$ 与 muon 衰变, 强调差别并比较电子能谱.
\end{exercise}
\begin{solution}
官方答案重编并补核子结构条件. 物理过程为 $n\to p e^-\bar\nu_e$; 原题省略反粒子横线. 两者都是三体带电弱衰变, 电子动能均连续. Muon 为基本轻子, neutron/proton 为复合核子, 后者矩阵元带 $V_{ud}$ 与 $g_A/g_V$, 不能只把 $m_\mu$ 换成 $m_n$. Muon 的电子总能量端点约 $m_\mu/2=52.8\,\mathrm{MeV}$; 中子的总能量端点约 $m_n-m_p=1.293\,\mathrm{MeV}$, 对应动能约 $0.782\,\mathrm{MeV}$.

Muon 谱为 $E_e^2(3m_\mu-4E_e)$, 无质量近似在端点附近达到最大; 中子允许谱为 $F(Z,E_e)p_eE_e(m_n-m_p-E_e)^2$, 在两端趋小, 最大值落在内部. 库仑畸变、反冲与核子结构修正控制精确形状和寿命.
\end{solution}
\begin{exercise}\label{ex:8701-4-3}
原题 3. 一般三阶及 n 阶幺正矩阵各有多少独立实参数? 三阶及 n 阶实正交矩阵呢? 可利用 $U=\e^{\ii H}$, H Hermitian.
\end{exercise}
\begin{solution}
官方解答重编. Hermitian 矩阵有 $n$ 个实对角元和 $n(n-1)/2$ 对复非对角元, 共 $n^2$ 个实参数, 三阶为 $9$. 实正交群单位连通分支由实反对称生成元产生, 共 $n(n-1)/2$ 个, 三阶为 $3$. 行列式 $\pm1$ 给两个离散分支, 不增加连续维数; 并非每个实正交矩阵都是实反对称矩阵的指数, 行列式负的分支需另乘一个反射.
\end{solution}
\begin{exercise}\label{ex:8701-4-4}
原题 4. LEP 在 $\sqrt s=M_Z$ 工作, 怎样不直接探测中微子衰变而由截面确定活跃中微子代数?
\end{exercise}
\begin{solution}
官方方法重编. 扫描共振曲线得到 $\Gamma_Z$ 和峰高, 可见轻子/强子道得到 $\Gamma_\ell,\Gamma_h$. 不可见宽度为 $\Gamma_Z-3\Gamma_\ell-\Gamma_h$, 于是
$N_\nu=(\Gamma_Z-3\Gamma_\ell-\Gamma_h)/\Gamma_\nu^{\rm SM}$.
峰强子截面 $12\pi\Gamma_e\Gamma_h/(M_Z^2\Gamma_Z^2)$ 提供独立关系. 单味 $\Gamma_\nu^{\rm SM}$ 必须用弱耦合计算; 原答案近似 $\Gamma_\nu\simeq2\Gamma_\ell$, 精密结果应保留 $\sin^2\theta_W$ 和辐射修正. 只计轻且能与 Z 耦合的态.
\end{solution}
\begin{exercise}\label{ex:8701-4-5}
原题 5. HERA 的电子、质子束流能量分别为 $27.5,820\,\mathrm{GeV}$. 用散射电子能量与角度表示 $Q^2,x,y$, 对 $3^\circ<\theta'_e<177^\circ$ 讨论覆盖区间.
\end{exercise}
\begin{solution}
官方公式重编及范围补充. 本章 $\theta'_e$ 相对入射电子方向, 质子沿 $+z$, 电子沿 $-z$. 忽略束流质量, $s=4E_eE_p=90200\,\mathrm{GeV}^2$,
\begin{equation}
 Q^2=4E_eE'_e\sin^2(\theta'_e/2),\quad
 y=1-\frac{E'_e}{E_e}\cos^2(\theta'_e/2),\quad x=\frac{Q^2}{sy}.
\end{equation}
固定 $E'_e$ 时角度端点给 $0.0754E'_e\lesssim Q^2\lesssim109.925E'_e$ (GeV$^2$, $E'_e$ 用 GeV). $x$ 用上式逐点计算, 不能把 $E'_e$ 未给定的函数范围误称固定数值边界.

物理区须 $0<y\le1,0<x\le1$, 给
$E'_e\le E_p/[\sin^2(\theta'/2)+(E_p/E_e)\cos^2(\theta'/2)]$.
只由角接受度不能给严格正的 $x_{\min}$、$Q^2_{\min}$: $E'_e\to0$ 时两者可趋零. 形式上 $x_{\max}=1$, 最大 $Q^2$ 在 $177^\circ$ 与该能量上界处约 $8.84\times10^4\,\mathrm{GeV}^2$. 实际探测电子能量阈值和 DIS 选择会进一步缩小区间, 原题未给这些条件. 当实验角度改用质子方向时应换为 $\pi-\theta'$.
\end{solution}
